\documentclass[10pt,a4paper]{amsart}
\usepackage[margin=19mm]{geometry}
\usepackage{amsmath,amssymb,mathtools}
\usepackage[scr=rsfs,cal=euler]{mathalfa}
\usepackage{xfrac}
\usepackage[unicode,hidelinks,bookmarks=false]{hyperref}
\newcommand{\ie}{i.e.\ }
\newcommand{\cf}{cf.\ }
\newcommand{\resp}{respectively\ }
\newcommand{\Subsection}{\subsection}
\providecommand{\nobreakdash}{\nobreak\hbox{-}}
\setlength{\emergencystretch}{2em}
\multlinegap0pt
\usepackage{mathtools}
\usepackage{amssymb}
\usepackage[shortlabels]{enumitem}
\usepackage{xcolor}
\usepackage{dsfont}
\usepackage{tikz}
\usepackage{tcolorbox}
\newtheorem{lemma}{Lemma}
\newtheorem{theorem}{Theorem}
\usepackage{cleveref}
\crefname{lemma}{Lemma}{Lemmas}
\crefname{theorem}{Theorem}{Theorems}
\DeclareMathOperator{\Cov}{Cov}
\newcommand{\R}{\mathbb{R}}
\DeclareMathOperator{\argmin}{argmin}
\DeclareMathOperator{\multiPPI}{MultiPPI}
\newcommand{\p}{\mathbb{P}}
\DeclareMathOperator{\supp}{supp}
\newcommand{\ul}[1]{\underline{#1}}
\newcommand{\wh}[1]{\widehat{#1}}
\title{Multiple-Prediction-Powered Inference}
\author{Charlie Cowen-Breen, Alekh Agarwal, Stephen Bates, William W. Cohen, Jacob Eisenstein, Amir Globerson, Adam Fisch}
\date{Source: arXiv:2603.27414v1 (CC BY 4.0)}
\begin{document}
\maketitle

\noindent\emph{Source and licence.} Multiple-Prediction-Powered Inference. Version arXiv:2603.27414v1 (CC BY 4.0), distributed by arXiv under the Creative Commons Attribution 4.0 International licence, http://creativecommons.org/licenses/by/4.0/, as recorded in the arXiv licence metadata for this exact version and shown on the abstract page https://arxiv.org/abs/2603.27414v1. Full licence notice: \texttt{licenses/arxiv-2603.27414-CC-BY-4.0.txt}.\par\smallskip

\noindent\textit{Editorial scope.} Counted unit: the article's theorem thm:asymptotic2 (source0.tex lines 3450--3455). The article's complete original proof of it is delivered with it (source0.tex lines 3458--3500, one \texttt{proof} environment of 43 lines, which the article places away from the statement). No mathematics is written, repaired or re-derived: the statement and the proof are the article's own lines, in the article's own notation, and no retained line is substituted or re-flowed.  Retained uncounted as the source-local context the proof needs: lines 381--384; lines 3322--3328; lines 3442--3447.  Nothing was omitted from any retained span, so the package carries no omission record.  No retained line contains a citation, so the wrapper carries no bibliography and no bibliography entry.  No retained line contains a figure environment, an image-inclusion command or drawing code, so no image is delivered and the package declares no asset dependency.  Editorial formatting only, in addition: an AMS \texttt{amsart} preamble that restates the article's own declarations and macros used by the retained lines, namely the environments \texttt{lemma}, \texttt{theorem} on separate counters, together with the standard packages those lines need.  The retained environments print in this document as Lemma 1, Theorem 1 in order of appearance, whereas the article prints the same items with the numbering of its own full document.  The complete native source of the article as distributed in the arXiv e-print for this exact version is bundled unchanged as \texttt{sources/197344/source0.tex}.
\begin{equation}
    \label{eq:variance}
    \mathcal{V}_B=\underset{\substack{\ul{n}\,:\,\text{\textbf{B} holds} \\ \supp(a)\subseteq\bigcup\{I:n_I>0\}}}{\min} a^\top S(\ul{n}) a,\quad\quad S(\ul{n})=\left(\sum_{I\in\mathcal{I}} n_I \Sigma_I^\dagger\right)^\dagger
\end{equation}

\begin{lemma}\label{lemma:discrete-convergence}
    Suppose that $\ul{\nu}^*(A^*)=\{\ul{{\nu}}^0\}$ is a singleton and that $A^{(N)}\to A^*$. Then for every $\epsilon>0$, there is some $N_1$ so that
    \[
    \left\| \frac{\ul{n}}{t}-\ul{\nu}^0 \right\| <\epsilon
    \]
    whenever $\ul{n}\in\ul{n}^*_t(A^{(N)})$ and $N,t\geq N_1$. 
\end{lemma}

\begin{equation}\label{eqn:condition}
\begin{split}
    \ul{\nu}^\star = &\argmin_{\ul{\nu}} V(\ul{\nu}) :=a^\top\left(\sum_I \nu_I P_I^\top \Sigma_I^{-1} P_I\right)^\dagger a \\
    & \text{s.t.}\quad \ul{\nu}\geq0,\quad \sum_I \nu_I c_I\leq B_0,\quad \supp(a)\subseteq\bigcup \{I:\nu_I>0\}
\end{split}
\end{equation}

\begin{theorem}[Asymptotic normality in the general setting]\label{thm:asymptotic2} Suppose $X\in\R^k$ has finite second moment, and suppose that $\Sigma=\Cov(X)$ satisfies \Cref{eqn:condition}. Suppose that $\wh{\Sigma}\overset{p}{\to}\Sigma$ in the operator norm as $t\to\infty$. Let $\hat{\theta}_{\multiPPI(\widehat{\Sigma})}$ denote the MultiPPI estimator with budget $tB_0$. Then for $\hat{\theta}_{\multiPPI(\widehat{\Sigma})}$ arbitrarily dependent on any potential samples used to estimate $\widehat{\Sigma}$, we have
    \[
    \sqrt{t}\left(\hat{\theta}_{\multiPPI(\widehat{\Sigma})}-\theta^*\right)\overset{d}{\to}\mathcal{N}(0,\mathcal{V}^*)
    \]
    as $B\to\infty$, where $\mathcal{V}^*=\lim_{B\to\infty} B\mathcal{V}_{B}$, and $\mathcal{V}_{B}$ is defined in \Cref{eq:variance}.
\end{theorem}

\begin{proof}[Proof of \Cref{thm:asymptotic2}]
% Let $\hat{\theta} = \hat{\theta}_{\text{MultiPPI}(\widehat{\Sigma})}$. Note that \ref{eqn:condition} is simply a rounded version of the optimization problem which is solved by $\hat{\theta}_{\text{MultiPPI}(\widehat{\Sigma})}$. Let $\wh{\ul{\nu}}$ denote the solution to \ref{eqn:condition} with $\Sigma$ replaced by $\wh{\Sigma}$.

% We first show that, as a result of the assumed condition, we have $\hat{\nu}_I\to \nu_I^*$ whenever $\wh{\Sigma}\to\Sigma$; that rounded solutions are optimal in the limit $t\to\infty$ is justified by \autoref{sec:rounding}. Since $a$ lies in the range of $\sum_I \nu^*_I P_I^\top \Sigma_I^{-1} P_I$, the objective function is continuous in $(\nu,\Sigma)$ at $\nu^*$. 

% The allocation $\underline{\hat{n}}$ and weights $\underline{\hat{\lambda}}$ are chosen to minimize the variance under $\widehat{\Sigma}$ subject to the budget $B = tB_0$. Let $\hat{\nu}_I = \hat{n}_I/t$. As $t \to \infty$, the optimal proportions $\underline{\hat{\nu}}$ converge to the solution $\underline{\nu}^*$ of the continuous optimization problem \ref{eqn:condition}. The convergence $\underline{\hat{\nu}} \xrightarrow{p} \underline{\nu}^*$ follows from $\widehat{\Sigma} \xrightarrow{p} \Sigma$ and Berge's Maximum Theorem, as the objective function is continuous and the feasible set is compact. By the continuous mapping theorem, we similarly have $\hat{\lambda}_I\overset{p}{\to} \lambda^*_I$.

% EDIT: 

For fixed $t$, we let allocation $\underline{\hat{n}}$ and weights $\underline{\hat{\lambda}}$ be chosen to minimize the variance of  $\hat{\theta} = \hat{\theta}_{\text{MultiPPI}(\widehat{\Sigma})}$ under $\widehat{\Sigma}$ subject to the budget $B = tB_0$, as in the MultiPPI procedure. Since $\wh{\Sigma}\overset{p}{\to}\Sigma$, \Cref{lemma:discrete-convergence} ensures that $\ul{\hat{n}}/t\overset{p}{\to}\ul{\nu}^*$, as defined in \Cref{eqn:condition}.

We can write $\sqrt{t}(\hat{\theta} - \theta^*) = \sum_{I \in \mathcal{I}} \hat{\lambda}_I^\top \sqrt{\frac{t}{\hat{n}_I}} W_{I, \hat{n}_I}$, where $W_{I, \hat{n}_I} = \frac{1}{\sqrt{\hat{n}_I}} \sum_{j=1}^{\hat{n}_I} (X_I^{(I,j)} - \mu_I)$. For indices $I$ with $\nu_I^* > 0$, we have $\hat{n}_I \xrightarrow{p} \infty$. Define $n_I^* = \lfloor t\nu_I^*\rfloor $, and let 
\[
W^*_I:=\frac{1}{\sqrt{n_I^*}}\sum_{i=1}^{n_I^*} (X_I^{(I,j)}-\mu_I)
\]
It is now enough to show that $W_{I,\hat{n}_I}-W^*_I\overset{p}{\to}0$, and this will follow from Kolmogorov's inequality. To simplify notation, let us focus on a single subset $I$, and define $Y_j = X_I^{(I,j)}-\mu_I$. Let us also define $S_m=\sum_{j=1}^m Y_j$. We must show that
\[
\frac{S_{\hat{n}}}{\sqrt{\hat{n}}}-\frac{S_{n^*}}{\sqrt{n^*}}\overset{p}{\to}0
\]
where we have dropped dependence on $I$ for convenience. We decompose
\[
\frac{S_{\hat{n}}}{\sqrt{\hat{n}}}-\frac{S_{n^*}}{\sqrt{n^*}} = \underbrace{\frac{S_{\hat{n}}-S_{n^*}}{\sqrt{n^*}}}_A + \underbrace{\frac{S_{\hat{n}}}{\sqrt{\hat{n}}}\left(1-\sqrt{\hat{n}/n^*}\right)}_B
\]
Fix $0<\delta<1$. We work on the event $E_\delta(t)=\{|\hat{n}-n^*|\leq \delta t\}$, which holds with high probability.

We first control $A$. On $E_\delta(t)$, $\sqrt{n^*}|A|$ is a sum of at most $\delta t+1$ i.i.d. copies of $Y_j$. Kolmogorov's inequality then yields
\[
\p\left(A>\epsilon \right) \leq \frac{\delta t+1}{\epsilon^2 n^*}\leq 4\frac{\delta}{\epsilon^2}
\]
because $n^*=\lfloor t\nu^*\rfloor$. Taking $\delta\to0$ yields that $A\overset{p}{\to}0$.

We next control $B$. Working again on $E_\delta(t)$, we have
\[
\frac{S_{\hat{n}}}{\sqrt{\hat{n}}} \leq \frac{1}{\sqrt{1-\delta}}\left(\underbrace{\frac{S_{n^*}}{\sqrt{n^*}}}_{O_p(1)}  + \underbrace{\frac{S_{\hat{n}}-S_{n^*}}{\sqrt{n^*}}}_A\right) 
\]
Recognizing the second term as $A\overset{p}{\to}0$, and the first term as tight by the central limit theorem, we conclude that $S_{\hat{n}}/\sqrt{\hat{n}}$ is tight. Now we conclude that $B\overset{p}{\to}0$ because $\hat{n}/n^*\overset{p}{\to}1$.

Having proven $W_{I,\hat{n}_I}-W^*_I\overset{p}{\to}0$, we conclude that
\[
\sqrt{t}(\hat{\theta}-\theta^*) = \sum_{I:\nu_I^*>0} \frac{1}{n^*_I} \sum_{j=1}^{n^*_I} (\lambda^*_I)^\top (X_I^{(I,j)}-\mu_I) + o_p(1)
\]
But this is precisely the desired result, since this is the solution to the continuous optimization problem, and we are done.
\end{proof}

\end{document}
